% ─────────────────────────────────────────────────────────────────────────
% §6 Limitations
% ─────────────────────────────────────────────────────────────────────────
\section{Limitations}\label{sec:limitations}

Five limitations are worth explicit disclosure.

\paragraph{Single-rater annotation.}
Labels are assigned by a single rater following the decision-tree
protocol described in §\ref{sec:method-labeling}. While 75.5\% of
responses were labeled without AI advisory and the full audit trail of
the remaining 24.5\% is released, this does not substitute for an
independent second rater. Readers should interpret all label-derived
quantities as rater-dependent. In §\ref{sec:method-reliability} we
explicitly decline to compute Cohen's $\kappa$ from the rater--judge
agreement on the AI-advised subset, as that subset is self-selected
(only borderline cases were submitted to the judge) and the judge was
trained on the same rulebook the rater consulted.
Inter-rater reliability with an independent rater is a priority for
replication.

\paragraph{Model-tier asymmetry.}
Our five-vendor panel is not capability-matched. \vendor{OpenAI} and
\vendor{Gemini} are represented by mini-class models
(\texttt{gpt-4o-mini}, \texttt{gemini-2.5-flash-lite}), while
\vendor{DeepSeek} and \vendor{Kimi} are represented by flagship models
(\texttt{deepseek-chat V3.2}, \texttt{kimi-k2-0905-preview}).
\vendor{Grok} sits at a middle tier (\texttt{grok-4-fast-non-reasoning}).
This asymmetry reflects our design choice to audit
vendor-as-recommended-production-endpoint rather than
capability-matched tiers, which we believe is a closer analog to the
typical downstream use case. We nonetheless cannot fully rule out that
some findings --- particularly the four-profile taxonomy of
§\ref{sec:results-finding5} --- partly reflect model-scale effects
rather than alignment-policy effects. A flagship-tier replication
subset is feasible at $\sim$USD 20 in API cost and represents a direct
path to disentangling these two factors; we defer it to a follow-up
paper.

\paragraph{Time-slice snapshot.}
Model identifiers were frozen on 2026-04-19 and are pinned in the
released source code. LLM vendor behavior evolves rapidly; an audit
conducted three or six months later may observe materially different
refusal patterns even with nominally the same model identifiers, and
vendors routinely update their safety stacks without public
announcement. Our findings describe the vendor behavior in the
April~2026 snapshot; longitudinal tracking is an explicit future-work
target.

\paragraph{Domain specificity.}
Our prompt bank covers Taiwan political sensitivity exclusively.
Generalization to other politically sensitive domains ---
e.g., Russian/Ukrainian, Palestinian/Israeli, or PRC mainland internal
political topics --- is not directly supported by our data. The
methodological contributions (the four-category refusal taxonomy, the
topic$\times$vendor$\times$layer decomposition, and the
paired-bootstrap inference protocol) are designed to be domain-portable
and we anticipate that replications on other domains will reveal a
similarly complex structure.

\paragraph{Modest prompt bank size.}
The 200-prompt bank is large enough for CI-disjoint inference on
aggregate-level vendor comparisons (Finding 1) but limited for
subgroup claims at small cell sizes. The Kimi API-blocked
analysis (§\ref{sec:results-finding2}) rests on 14 prompts with
per-topic subcounts ranging from 0 to 9. The four-profile taxonomy
(§\ref{sec:results-finding5}) relies on sovereignty-topic cells with
39 prompts per vendor; bootstrap CIs on sovereignty-specific on-task
rates are correspondingly wide (up to 31 percentage-points for
\vendor{DeepSeek}'s \ci{2.6}{23.3}). All findings argued from
subgroup analyses report CIs explicitly and do not rely on
point-estimate claims; we recommend a scaled-up replication on
$N \geq 500$ prompts for any follow-up that seeks to characterize
sub-topic structure with higher resolution.

% ─────────────────────────────────────────────────────────────────────────
% §7 Conclusion and Future Work
% ─────────────────────────────────────────────────────────────────────────
\section{Conclusion and Future Work}\label{sec:conclusion}

We have presented an audit of refusal behavior across five commercial
large language models on a fixed bank of 200 Taiwan-political prompts.
The audit's methodological design --- prompt-level paired bootstrap,
four-class response taxonomy including an API-level filter category,
and topic$\times$vendor$\times$layer stratification --- was designed
to make aggregate-level and mechanism-level vendor comparisons directly
comparable, and the resulting findings consistently favor a
multi-dimensional over a single-axis interpretation of alignment
culture.

Three substantive conclusions follow. First, vendor refusal behavior
does not partition along geographic (East--West) or regulatory lines in
our data; the two Chinese-owned vendors in the panel produce the most
distant pair of refusal distributions in the matrix, and the aggregate
distribution closest to U.S.\ vendors belongs to \vendor{DeepSeek}.
Second, vendor alignment strategies span two architecturally distinct
layers (infrastructure filtering and model-level RLHF), with
\vendor{Kimi} the single vendor in our panel operating both at scale;
the two layers are not commensurable as refusal-rate contributors and
should be reported separately in future audits. Third, a
topic-stratified view reveals structure invisible to aggregate analysis:
\vendor{DeepSeek} exhibits a sovereignty-specific on-task collapse
(10.3\% vs 54.0\% non-sovereignty, disjoint-CI evidence) that has no
analog in any other vendor in our panel.

Four future-work targets suggest themselves. \textbf{First}, independent
re-annotation with multiple raters to generate inter-rater reliability
statistics; the released decision-tree rulebook and AI-judge audit
trail are intended to facilitate this. \textbf{Second}, a capability-matched
replication in which each vendor is queried at its flagship-tier model
rather than at its recommended production endpoint, to disentangle
alignment-policy from model-scale effects on the observed findings.
\textbf{Third}, longitudinal replication across model revisions to
characterize the stability of the four-profile taxonomy. \textbf{Fourth},
domain replication --- particularly to PRC-internal political topics,
where the Taiwan-statehood characterization of \vendor{Kimi}'s filter
would predict a different trigger pattern --- to test the generality of
the topic$\times$vendor$\times$layer framework.

The project's broader motivation was the construction of LLM agent
populations for Taiwan voter simulation; the present audit quantifies
the vendor-dependent priors against which any such simulation study
should be interpreted. We recommend that downstream agent-simulation
work in politically sensitive domains treat vendor as a first-class
experimental variable, run multi-vendor replications where feasible,
and report layer-stratified refusal metrics.
